\section{Special Matrix Ring: Identity and Inverse Computations}

\begin{example}[A Commutative Matrix Ring with Non-Standard Unity]
Consider the set:
\[
A = \left\{ \begin{bmatrix} x & x \\ x & x \end{bmatrix} \;\middle|\; x \in \mathbb{R} \right\}
\]
under standard matrix addition and multiplication.
\end{example}

\subsection{Verification of Ring Axioms}

\begin{enumerate}
    \item \textbf{$(A, +)$ is an abelian group:}
    For any $X = \begin{bmatrix} x & x \\ x & x \end{bmatrix}$ and $Y = \begin{bmatrix} y & y \\ y & y \end{bmatrix} \in A$:
    \[
    X + Y = \begin{bmatrix} x+y & x+y \\ x+y & x+y \end{bmatrix} \in A.
    \]
    Addition is associative and commutative since addition in $\mathbb{R}$ is. The zero element is:
    \[
    \mathbf{0} = \begin{bmatrix} 0 & 0 \\ 0 & 0 \end{bmatrix} \in A,
    \]
    and the additive inverse of $X$ is $-X = \begin{bmatrix} -x & -x \\ -x & -x \end{bmatrix} \in A$.
    
    \item \textbf{Commutativity of multiplication in $A$:}
    Let $X, Y \in A$ with $X = \begin{bmatrix} x & x \\ x & x \end{bmatrix}$ and $Y = \begin{bmatrix} y & y \\ y & y \end{bmatrix}$. Then:
    \begin{align*}
    \begin{bmatrix} x & x \\ x & x \end{bmatrix} \begin{bmatrix} y & y \\ y & y \end{bmatrix} &= \begin{bmatrix} x \cdot y + x \cdot y & x \cdot y + x \cdot y \\ x \cdot y + x \cdot y & x \cdot y + x \cdot y \end{bmatrix} \\
    &= \begin{bmatrix} 2xy & 2xy \\ 2xy & 2xy \end{bmatrix} \\
    &= \begin{bmatrix} 2yx & 2yx \\ 2yx & 2yx \end{bmatrix} \\
    &= \begin{bmatrix} y & y \\ y & y \end{bmatrix} \begin{bmatrix} x & x \\ x & x \end{bmatrix}.
    \end{align*}
    Therefore, multiplication is commutative in $A$.
    
    \item \textbf{Distributivity:}
    Matrix multiplication is always distributive over matrix addition.
\end{enumerate}
Thus, $(A, +, \cdot)$ is a \textbf{commutative ring}.

\subsection{Finding the Unity Element}

Let $E \in A$ be the multiplicative identity (unity) of $A$, so that:
\[
E = \begin{bmatrix} e & e \\ e & e \end{bmatrix}.
\]
Then $E$ must satisfy $X E = X$ for all $X \in A$:
\[
\begin{bmatrix} x & x \\ x & x \end{bmatrix} \begin{bmatrix} e & e \\ e & e \end{bmatrix} = \begin{bmatrix} x & x \\ x & x \end{bmatrix}.
\]
Performing the multiplication:
\[
\begin{bmatrix} 2xe & 2xe \\ 2xe & 2xe \end{bmatrix} = \begin{bmatrix} x & x \\ x & x \end{bmatrix}.
\]
Equating corresponding entries:
\[
2xe = x \implies 2e = 1 \implies e = \frac{1}{2}.
\]
Therefore, the identity element of the ring $A$ is:
\[
E = \begin{bmatrix} 1/2 & 1/2 \\ 1/2 & 1/2 \end{bmatrix}.
\]
\begin{remark}
Notice that the unity element $E \in A$ is \emph{not} the standard $2 \times 2$ identity matrix $I_2 = \begin{bmatrix} 1 & 0 \\ 0 & 1 \end{bmatrix}$ (which does not belong to $A$), yet it functions as the exact multiplicative identity on $A$.
\end{remark}

\subsection{Computation of Multiplicative Inverses}

For $(A \setminus \{\mathbf{0}\}, \cdot)$ to form an abelian group (making $A$ a field), each non-zero element must possess a multiplicative inverse. Let:
\[
Y = X^{-1} = \begin{bmatrix} y & y \\ y & y \end{bmatrix}.
\]
Then $X Y = E$:
\[
\begin{bmatrix} x & x \\ x & x \end{bmatrix} \begin{bmatrix} y & y \\ y & y \end{bmatrix} = \begin{bmatrix} 1/2 & 1/2 \\ 1/2 & 1/2 \end{bmatrix}.
\]
This gives:
\[
\begin{bmatrix} 2xy & 2xy \\ 2xy & 2xy \end{bmatrix} = \begin{bmatrix} 1/2 & 1/2 \\ 1/2 & 1/2 \end{bmatrix},
\]
which implies:
\[
2xy = \frac{1}{2} \implies y = \frac{1}{4x} \quad (x \neq 0).
\]
Therefore, for any non-zero element $X \in A$, its multiplicative inverse is:
\[
X^{-1} = \begin{bmatrix} \dfrac{1}{4x} & \dfrac{1}{4x} \\[6pt] \dfrac{1}{4x} & \dfrac{1}{4x} \end{bmatrix}.
\]
Consequently, $(A \setminus \{\mathbf{0}\}, \cdot)$ is an abelian group, and $A$ is isomorphic as a field to $\mathbb{R}$.
